\section{Contesto}\label{sec:background}

\subsection{Il Sistema di Registrazione AIRE e i suoi Limiti Strutturali}
L'\emph{Anagrafe degli Italiani Residenti all'Estero} è stata istituita dalla Legge 470/1988~\cite{legge470_1988} per mantenere un registro dei cittadini italiani che vivono all'estero per più di dodici mesi consecutivi.
L'iscrizione è un obbligo di legge: i cittadini devono informare il proprio comune di residenza prima di partire o, in alternativa, contattare il consolato italiano più vicino entro novanta giorni dallo stabilimento della residenza all'estero.
I registri comunali alimentano le statistiche ufficiali sull'emigrazione dell'ISTAT, rendendo l'AIRE la fonte primaria per qualsiasi studio quantitativo sull'emigrazione italiana~\cite{istat_migrazioni2023}.
In linea di principio, il sistema è completo.
In pratica, soffre di un divario di conformità ampiamente documentato.
L'iscrizione richiede iniziativa da parte dell'emigrante: non conferisce alcun beneficio amministrativo immediato nel paese di destinazione e la mancata iscrizione non comporta sanzioni significative.
Il risultato è un classico problema di auto-segnalazione, descritto nella letteratura sulla qualità dei dati come \emph{sottostima amministrativa}: il registro cattura solo quegli individui che hanno scelto, o si sono ricordati, di interagire con una procedura burocratica percepita come facoltativa.
La Fondazione Migrantes, che pubblica un rapporto annuale sui dati dell'emigrazione italiana, ha costantemente riferito che i totali AIRE sono inferiori alle stime derivate dai registri consolari e dalle fonti dei paesi di destinazione~\cite{migrantes2022}.
\subsection{Le Statistiche Speculari come Strumento di Misurazione}
Quando una singola fonte amministrativa è nota per essere incompleta, un approccio standard nella contabilità demografica è quello di convalidarla incrociandola con una fonte indipendente che registri lo stesso evento dal lato opposto.
Questo metodo è stato applicato a livello europeo da Eurostat, che riconosce esplicitamente come i flussi di emigrazione siano più difficili da misurare rispetto ai flussi di immigrazione: la partenza viene raramente registrata con la stessa urgenza amministrativa dell'arrivo.
Questo perché le istituzioni del paese ricevente — permessi di soggiorno, registrazione fiscale, iscrizione scolastica — creano incentivi naturali affinché gli immigrati si rendano visibili alle autorità, incentivi che non esistono sul lato della partenza~\cite{eurostat_migration_explained2023, ec_regulation_862_2007}.
Questa asimmetria è strutturale, ed è esattamente ciò che sfrutta la nostra analisi.
